problem:  
Let \((M,\Delta,g)\) be a compact, connected, equiregular **step‑2 sub‑Riemannian manifold** of dimension \(n\) and rank \(r<n\). Let  
\[
H(x,p) = \tfrac{1}{2}\langle p|_{\Delta_x},p|_{\Delta_x}\rangle_g
\]
be the sub‑Riemannian Hamiltonian on \(T^*M\), and \(\phi_t^H\) its Hamiltonian flow.

Suppose that there exists a smooth **Abelian Lie algebra \(\mathfrak{a}\)** of complete vector fields on \(M\) of dimension \(k\geq 1\) satisfying:

1. each \(X\in \mathfrak{a}\) preserves the sub‑Riemannian structure:  
   \([X,\Gamma(\Delta)]\subset \Gamma(\Delta)\) and \(L_X g=0\) on \(\Delta\);

2. the corresponding **momentum functions** \(I_X(x,p) = \langle p, X(x)\rangle\) on \(T^*M\) Poisson‑commute with \(H\): \(\{H,I_X\}=0\) for every \(X\in \mathfrak{a}\);

3. the lifted flows of \(\mathfrak{a}\) act freely and symplectically on some regular energy level \(\Sigma=H^{-1}(\tfrac12)\), defining a smooth symplectic quotient space \(\Sigma/\mathfrak{a}\).

Conjectural question:  
If \(k\) is maximal with these properties and the reduced Hamiltonian system on \(\Sigma/\mathfrak{a}\) is **Liouville‑integrable** in the classical sense (i.e., admits \(n-k=\frac12\dim(\Sigma/\mathfrak{a})\) smooth independent commuting integrals), must \((M,\Delta,g)\) locally arise as the **horizontal geometry of a Riemannian submersion**  
\[
\pi:(M,g_R)\to (B,g_B),
\]
where \(g_R\) is a Riemannian metric whose geodesic Hamiltonian flow projects to \(\phi_t^H\) on horizontal covectors?  
Equivalently, does complete integrability generated by linear (momentum) first integrals and their reduction force the sub‑Riemannian Hamiltonian to be locally that of a Riemannian submersion metric?

Why is it a "good" problem:  
This problem is both rigorous and open, representing a refined version of the theme “integrability rigidifies geometry.”  

- **Logical soundness:** The integrability framework now uses a clearly defined **group‑action reduction**: the linear integrals come from true symmetries on \(M\), their momentum map provides a Hamiltonian group action, and Liouville integrability is imposed only after proper symplectic reduction.  
- **Profound insight:** It probes whether *enough symmetries and residual integrability* of the sub‑Riemannian geodesic flow imply that the supposed “nonholonomic” geometry is in fact the shadow of a Riemannian submersion. This connects dynamical integrability with intrinsic geometric rigidity.  
- **Bridge between fields:** The question unites **Hamiltonian dynamics and symplectic reduction**, **sub‑Riemannian and Riemannian geometry**, and **differential topology of foliations**. It brings techniques from geometric control and integrable systems into contact with the theory of Riemannian submersions and foliation geometry.  
- **Simplicity with depth:** The formulation—requiring commuting Killing symmetries and classical Liouville integrability on the reduced system—is conceptually concise yet leads to a high‑dimensional, intricate rigidity phenomenon that remains unexplored.  
- **Potential impact:** A positive resolution would characterize the only Liouville‑integrable sub‑Riemannian systems with homogeneous linear integrals as those derived from Riemannian submersions, uniting symmetry, integrability, and metric geometry in a single rigidity principle.